\documentclass[10pt,a4paper]{amsart}
\usepackage[margin=19mm]{geometry}
\usepackage{amsmath,amssymb,mathtools}
\usepackage[scr=rsfs,cal=euler]{mathalfa}
\usepackage{xfrac}
\usepackage[unicode,hidelinks,bookmarks=false]{hyperref}
\newcommand{\ie}{i.e.\ }
\newcommand{\cf}{cf.\ }
\newcommand{\resp}{respectively\ }
\newcommand{\Subsection}{\subsection}
\providecommand{\nobreakdash}{\nobreak\hbox{-}}
\setlength{\emergencystretch}{2em}
\multlinegap0pt
\usepackage{amssymb}
\usepackage{xcolor}
\usepackage{tikz}
\usepackage{enumerate}
\usepackage[shortlabels]{enumitem}
\newtheorem{lemma}{Lemma}
\newcommand{\C}{\mathbb{C}}
\DeclareMathOperator{\Prim}{Prim}
\newcommand{\TT}{\mathbb{T}}
\DeclareMathOperator{\U}{U}
\newcommand{\Z}{\mathbb{Z}}
\newcommand{\ZZ}{\mathcal{Z}}
\DeclareMathOperator{\diag}{diag}
\DeclareMathOperator{\ind}{ind}
\newcommand{\mt}{\mathtt}
\title{Nuclear Dimension of Twisted $C^*$-Algebras of Virtually Abelian Groups}
\author{Forrest Glebe, Pradyut Karmakar, Iason Moutzouris}
\date{Source: arXiv:2605.27936v2 (CC BY 4.0)}
\begin{document}
\maketitle

\noindent\emph{Source and licence.} Nuclear Dimension of Twisted $C^*$-Algebras of Virtually Abelian Groups. Version arXiv:2605.27936v2 (CC BY 4.0), distributed by arXiv under the Creative Commons Attribution 4.0 International licence, http://creativecommons.org/licenses/by/4.0/, as recorded in the arXiv licence metadata for this exact version and shown on the abstract page https://arxiv.org/abs/2605.27936v2. Full licence notice: \texttt{licenses/arxiv-2605.27936-CC-BY-4.0.txt}.\par\smallskip

\noindent\textit{Editorial scope.} Counted unit: the article's lemma group (source0.tex lines 653--661). The article's complete original proof of it is delivered with it (source0.tex lines 687--689, one \texttt{proof} environment of 3 lines, which the article places away from the statement). No mathematics is written, repaired or re-derived: the statement and the proof are the article's own lines, in the article's own notation, and no retained line is substituted or re-flowed.  No further source-local context is needed: every cross-reference of every retained line resolves inside the two delivered spans.  Nothing was omitted from any retained span, so the package carries no omission record.  No retained line contains a citation, so the wrapper carries no bibliography and no bibliography entry.  No retained line contains a figure environment, an image-inclusion command or drawing code, so no image is delivered and the package declares no asset dependency.  Editorial formatting only, in addition: an AMS \texttt{amsart} preamble that restates the article's own declarations and macros used by the retained lines, namely the environments \texttt{lemma} on separate counters, together with the standard packages those lines need.  The retained environments print in this document as Lemma 1 in order of appearance, whereas the article prints the same items with the numbering of its own full document.  The complete native source of the article as distributed in the arXiv e-print for this exact version is bundled unchanged as \texttt{sources/201508/source0.tex}.
\begin{lemma}\label{existance of irrep on finite group}
  There exists $m\in \Z$ and $\pi\in \Prim_m(C^*(F))$ such that $(\pi \circ q)(a)=\diag(\omega,\ldots,\omega)$.
  \begin{proof}
      Recall that $q$ is the natural map from $L$ to $F$. Start with the character $\pi_0 \colon \langle q(a)\rangle\to \TT$ defined via $\pi_0(q(a))=\omega$. Set $\pi_1:=\ind_{\langle q(a)\rangle}^F \pi_0:L\to \U(|F|/n)$. By construction of the induced representation, and the fact that $a\in \ZZ(L)$, we deduce that $\pi_1(q(a))=\diag(\omega,...,\omega)$. \par

      
      Notice that $\pi_1$ is a direct sum of irreducible representations. Let $\pi \colon F\to U_m(\C)$ be one of these. The proof is complete because $\pi(q(a))=\diag(\omega,...,\omega)$.
  \end{proof}
\end{lemma}

    \begin{proof}
        It is clear from the above construction that $\Phi(U)\subset \Prim_{\mt{K}\cdot m}(C^*(G))$. Because $a$ has finite order, it follows that $\iota(a)=(e,b)$ for some $b\in F$ such that $q(a)=b$. Thus, for every $\chi\in U$, we have $\varphi_{\chi}\circ \iota(a)=\pi(b)=(\pi\circ q)(a)=\diag(\omega,\omega,...,\omega)$. Note that the last equality follows from Lemma \ref{existance of irrep on finite group}. 
    \end{proof}

\end{document}
